% Résumé en cinq mouvements (trame, étape 11) : contexte, faille, résultat majeur, résultats-clés, portée.
\begin{abstract}
\noindent Les agents de recherche entraînés par renforcement, comme Search-R1, répondent mieux que les systèmes à règle fixe, et l'industrie parie sur ce type d'entraînement. Mais on ne sait pas ce qu'il apprend : rend-il possibles des réponses que le modèle de départ ne trouvait jamais, ou rend-il seulement plus fiables celles qu'il trouvait déjà ? Les études existantes portent sur des agents entraînés par leurs auteurs, sans test apparié. Nous comparons deux agents Search-R1 publics à leur propre modèle de départ, Qwen2.5-3B-Instruct, sur 500 questions HotpotQA à plusieurs étapes, avec 32 essais par question et un plan d'analyse gelé avant les runs. Le renforcement fiabilise sans élargir : l'agent gagne 4 à 7 points au premier essai, mais à 32 essais il ne résout pas plus de questions que son modèle de départ (351 et 350 contre 346). Il apprend surtout à chercher : il ramène les deux paragraphes utiles dans 54~\% des cas, contre 32~\%, et ce gain explique toute son avance. Il lit aussi mieux quand on l'oblige à répondre, mais renonce à répondre dans 14 à 20~\% des cas où il a l'information, ce qui annule ce gain. Sur des questions de comparaison, il ne gagne rien. Une première étude, à modèles de départ différents, attribuait l'avantage à la lecture ; la comparaison exacte corrige cette conclusion. Si le renforcement d'un agent de recherche fiabilise surtout sa recherche, un entraînement supervisé sur ses bonnes traces pourrait en obtenir l'essentiel à moindre coût.
\end{abstract}

\begin{otherlanguage}{english}
\begin{abstract}
\noindent Search agents trained with reinforcement learning (RL), such as Search-R1, outperform fixed-rule systems, and industry is betting on this kind of training. Yet what it teaches remains unclear: does it make answerable what the base model never answered, or only make existing answers more reliable? Prior studies use agents trained by their own authors, without paired tests. We compare two public Search-R1 agents with their own base model, Qwen2.5-3B-Instruct, on 500 multi-hop HotpotQA questions, with 32 attempts per question, a closed search engine and an analysis plan frozen before the runs. RL sharpens without expanding: the agent gains 4 to 7 points at the first attempt, but with 32 attempts it solves no more questions than its base model (351 and 350 versus 346). It mostly learns to search: it retrieves both supporting paragraphs in 54\% of trajectories versus 32\%, which accounts for its whole advantage. It also reads better when forced to answer, but gives up answering in 14 to 20\% of the cases where it has the information, which cancels that gain. On comparison questions it gains nothing. A first study, with different base models, credited reading; the exact comparison corrects it. If RL on a search agent mainly makes its search more reliable, supervised training on its successful traces might capture most of the benefit at a lower cost.
\end{abstract}
\end{otherlanguage}
